% -----------------------------------------------------------
\section{Mercy}
% -----------------------------------------------------------
	\label{MERCY}
	
% % ----------------------
% \subsection{See also}
% % ----------------------

% % ----------------------
% \subsection{1828 American Dictionary of the English Language} % *1828
% % ----------------------
	% \label{MERCY-1828}

	% \begin{compactitem}
		% \item
	% \end{compactitem}

% % ----------------------
% \subsection{Oxford English Dictionary} % *OED
% % ----------------------
	% \label{MERCY-OED}

	% \begin{compactitem}
		% \item[]
	% \end{compactitem}

% ----------------------
\subsection{Scriptural Usage} % *SCRIPTURES
% ----------------------
	\label{MERCY-SCRIPTURES}

	% -----
	\subsubsection{Old Testament} % *OT
	% -----
		\label{MERCY-OT}
	
		% % --
		% \paragraph{KJV}
		% % --
			% \label{MERCY-OT-KJV}
	
			% \begin{tabular}{ll}
				% Total occurrences in the OT & 
			% \end{tabular}
			
			% \begin{compactdesc}
				% \item[]
			% \end{compactdesc}
			
		% --
		\paragraph{Hebrew Bible}
		% --
			\label{MERCY-HEBREW}
		
			%
			\subparagraph{\texorpdfstring{\he{.hEsEd}}{}} % חֶסֶד
			%
				\label{CHESED} % whatever is in step bible without periods
			
				\begin{compactdesc}
					\item[HALOT 336b, PDF 393] \phantom{.}
						\begin{compactitem}
							\item[1] \textbf{joint obligation} between relatives, friends, host and guest, mater and servant; closeness, solidarity, \textbf{loyalty}
							\item[1.a] \he{.hEsEd} and \he{b*:riyt}; \he{b*:riyt} comes about by a ceremony \he{.hEsEd} results from the closer relationship between two people, the obligations are largely the same...lasting loyalty, \textbf{faithfulness}
							\item[1.b] \he{.hEsEd} exists between a son and a dying father \hyperref[GENESIS-47-29]{Gen 47:29}, a wife and a husband \hyperref[GENESIS-20-13]{Gn 20:13}, relatives \hyperref[RUTH-2-20]{Ru 2:20}...[and others]
							\item[2] \he{.hEsEd} in God's relationship with the people or an individual, \textbf{faithfulness, goodness, graciousness}
							\item[2.a] \he{.hEsEd yhwh}
							\item[2.b] \he{`A,sAh .hEsEd} to show faithfulness...
							\item[3] plural \he{.ha:sAdiyM}, \he{.ha:sAday} etc. the individual actions resulting from solidarity
							\item[3.a] (of people) \textbf{godly action, achievements}
							\item[3.b] (God's) \textbf{proofs of mercy}
						\end{compactitem}
					\item[BDB 338b, PDF 357] \phantom{.}
						\begin{compactitem}
							\item goodness, kindness
							\item[I] \textit{of man}
							\item[I.1] \textit{kindness} of men towards men, in doing favours and benefits
							\item[I.2] \textit{kindness} (especialy as extended to the lowly, needy and miserable), \textit{mercy}
							\item[I.3] (rarely) \textit{affection} of Israel to \he{yhwh}, \textit{love to God, piety}
							\item[I.4] \textit{lovely appearance}
							\item[II] \textit{of God: kindness, lovingkindness} in condescending to the needs of his creatures. He is \he{.has:d*AM} [or] \textit{their goodness, favour} \hyperref[JONAH-2-9]{Jon 2:9}...
							\item[II.1] specifically lovingkindness
							\item[II.1.a] \textit{in redemption from enemies and troubles}
							\item[II.1.b] \textit{in preservation of life from death}
							\item[II.1.c] \textit{in quickening of spiritual life}
							\item[II.1.d] \textit{in redemption from sin}
							\item[II.1.e] \textit{in keeping the covenants}
							\item[II.2] \he{.hEsEd} is grouped with other divine attributes: \he{.hsd w'mt} \textit{kindness (loving-kindness) and fidelity} \hyperref[GENESIS-24-27]{Gn 24:27}...
							\item[II.3] the \textit{kindness} of God is
							\item[II.3.a] \textit{abundant:} \he{rab--.hEsEd} \textit{abundant, plenteous in kindness (goodness)}...it is kept for thousands \hyperref[EXODUS-34-7]{Ex 34:7}...for 1000 generations \hyperref[DEUTERONOMY-7-9]{Dt 7:9}...
							\item[II.3.b] \textit{great in extent}
							\item[II.3.c] \textit{everlasting}
							\item[II.3.d] \textit{good}
							\item[II.4] plural \textit{mercies, deeds of kindness}, the historic displays of loving-kindness to Israel
						\end{compactitem}
					\item[DCH PDF 1412] \phantom{.}
						\begin{compactitem}
							\item \textbf{loyalty}---\textbf{loyalty, faithfulness, kindness, love, mercy}, plural \textbf{mercies, (deeds of) kindness}, of Y. to humans...of humans to Y...between humans...of things.
						\end{compactitem}
					\item[TDOT] (See my annotated PDF \href{https://drive.google.com/file/d/1oe9t0mj7ZDi8Yf0bM4C4g2NB0uM5sfIH/view?usp=share_link}{here}.)
					\item[TLOT] \phantom{.}
						\begin{compactdesc}
							\item[PDF 596] The word \he{.hEsEd}...is only insufficiently rendered by the Eng. term ``kindness.''
							\item[PDF 598, 600] According to Glueck, \he{.hEsEd} does not refer to a spontaneous, ultimately unmotivated kindness, but to a mode of behavior that arises from a relationship defined by rights and obligations (husband-wife, parent-child, prince-subjects). When \he{.hEsEd} is attributed to God, it concerns the realization of the promises inherent in the covenant. When \he{.hEsEd} does assume connotations of kindness, it is the result of a secondary assimilation to \he{rA.hamiyM}...Furthermore, this view means that the formulation \he{.hEsEd wE'E:mEt} should be construed as a hendiadys...To the extent, then, that one can speak of a blurring of the linguistic boundaries between \he{.hEsEd} and \he{rA.hamiyM}, it seems to have resulted in more of a limitation than an expansion of meaning.
							\item[PDF 601] The observations offered to this point concerning the occurrence of the term, its development and delimitation, and elucidating supplementations of it allow the conclusion, for the present still indefinite, that \he{.hEsEd} means some special reciprocal behavior, something that exceeds the matter of course. This conclusion may be demonstrated and substantively illustrated in detail by a survey of passages, esp. in narrative literature, in which \he{.hEsEd} is exercised in the interpersonal realm.
							\item[PDF 604] It is not possible to convey precisely \he{.hEsEd}'s semantic range as encountered in profane usage with one Eng. word. \he{.hEsEd} is not ``grace,'' and the often suggested ``favor'' is insufficient. \hl{First, }\he{.hEsEd} \hl{occurs tangibly in concrete situations and at the same time transcends the individual demonstration and envisions the actor.} In this regard, the term exhibits affinities with Eng. ``kindness,'' and with ``goodness'' as well (see 3c). \he{.hEsEd} does occur in relation to particular social forms; its configuration may even be governed by them, but it is never the obvious, the obligatory. \hl{It is a human demeanor that alone can fill a form with life and is in some cases (not always) the very requirement for the birth of a community. Jepsen (op. cit. 269) attempted to describe the intention as good will that becomes good deed. This notion is certainly included but is alone insufficient. I suggest an expression for magnanimity, for a sacrificial, humane willingness to be there for the other} (Stoebe, diss. 67; id., \textit{VT} 2 [1952]: 248). \hl{It is a given that} \he{.hEsEd} \hl{always has to do in some way with the life of the other, and that one expects and hopes from the recipient of such} \he{.hEsEd} \hl{a similar willingness that in turn surpasses the obligatory.}
							\item[PDF 607--608] Concerning the question of the reciprocity of \he{.hEsEd}, one should first of all refer to \hyperref[HOSEA-10-12]{Hos 10:12}: ``Sow righteousness and harvest \he{.hEsEd}.'' \he{.s:dAqAh} and \he{.hEsEd} will be given by God, on the one hand, and are, on the other, to be actualized by people, so that God's \he{.hEsEd} is both requirement and example of the proper human attitude toward him...

								\hyperref[HOSEA-12-7]{Hos 12:7} belongs here too. The demand ``practice \he{.hEsEd} and justice'' is included in a call to return to God and should be understood in this context. Although the emphasis here seems to lie more pronouncedly on the behavior of people toward each other, the two are still not to be strictly distinguished. \hyperref[HOSEA-6-6]{Hos 6:6} also demonstrates this duality in its juxtaposition of \he{.hEsEd} and sacrifice (cf. \hyperref[1SAMUEL-15-22]{1 Sam 15:22}). The alternative of \he{.hEsEd} toward God or only among people is falsely posed because, for the OT, both belong together. That \he{.hEsEd} and the sacrificial cult are so juxtaposed should be understood against the background of the fact that sacrifice need not exclude human devotion, but that it can also be understood as a duty with necessary consequences for behavior toward others too (cf. \hyperref[AMOS-8-4]{Amos 8:4--6}). 

								Esp. interesting in this respect is \hyperref[HOSEA-4-1]{Hos 4:1}. The reversal of the sequence in which the terms from \hyperref[HOSEA-2-21]{2:21f.} appear in the demand directed at people is intentional: they form a descending climax. Yahweh contends on his own behalf: if the people have no constancy, they should at least have \he{.hEsEd} devotion; if this, too, is lacking, they should at least have an awareness of what Yahweh has done and given. \hyperref[HOSEA-6-4]{Hos 6:4} illustrates that these thoughts are not otherwise unknown to Hosea. The influence of God's punitive act engenders something like a \he{.hEsEd} attitude (contra Jepsen, op. cit. 269). One at least takes him into account. But this attitude has as little duration as dew or morning fog. 
						\end{compactdesc}
					\item[TWOT] \phantom{.}
						\begin{compactdesc}
							\item[305--306] For centuries the word \he{.hEsEd} was translated with words like mercy, kindness, love. The LXX usually uses \textit{eleos} ``mercy,'' and the Latin \textit{misericordia}. The Targum and Syriac use frequently a cognate of \he{.tOb}. The root is not found in Akkadian or Ugaritic. The lexicons up through BDB and GB (which said \textit{Liebe, Gunst, Gnade}, love, goodness, grace) are similar. KB however is the ``mutual liability of those...belonging together.''

								In 1927 Nelson Glueck, shortly preceded by I. Elbogen, published a doctoral dissertation in German translated into English by A. Gottschalk, \textit{Hesed in the Bible} with an introduction by G. A. LaRue which is a watershed in the discussion. His views have been widely accepted. In brief, Glueck built on the growing idea that Israel was bound to its deity by covenants like the Hittite and other treaties. He held that God is pictured as dealing basically in this way with Israel. The Ten Commandments, etc. were stipulations of the covenant, lsrael's victories were rewards of covenant keeping, her apostasy was covenant violation and God's \he{.hEsEd} was not basically mercy, but loyalty to his covenant obligations, a loyalty which the lsraelites should also show. He was followed substantially by W. F. Lofthouse (1933), N. H. Snaith (1944), H. W. Robinson (1946), Ugo Masing (1954), and many others.

								There were others, however, who disagreed. F. Assension (1949) argued for mercy, basing his views on the \textsc{OT} versions. H. J. Stoebe (doctoral dissertation 1951, also articles in 1952 VT and in THAT) argued for good-heartedness, kindness [this is the TLOT entry above, Stoebe is the author]. Sidney Hills and also Katherine D. Sakenfeld...held in general that \he{.hEsEd} denotes free acts of rescue or deliverance which in prophetic usage includes faithfulness...

								\hl{The writer [of this entry in TWOT] would stress that the theological difference is considerable whether the Ten Commandments are stipulations to a covenant restricted to Israel to which God remains true and to which he demands loyalty, or whether they are eternal principles stemming from God's nature and his creation to which all men are obligated and according to which God will judge in justice or beyond that will show love, mercy and kindness}...

								Glueck makes something of \hyperref[1SAMUEL-20-8]{I Sam 20:8, 14, 15} where David and Jonathan swore friendship. This covenant, says G[lueck]. was the basis of the \he{.hEsEd}. Here, perhaps, is G's major mistake. \hl{He forgets that covenants arise on the basis of a relationship and that the obligations are often deeper than the covenant.}

							\item[307] \hl{So, it is obvious that God was in covenant relation with Israel, also that he expressed this relation in} \he{.hEsEd}\hl{, that God's} \he{.hEsEd} \hl{was eternal (Note the refrain of} \hyperref[PSALMS-136]{Ps 136})\hl{---though the }\he{.hEsEd} \hl{of Ephraim and others was not (}\hyperref[HOSEA-6-4]{Hos 6:4}\hl{). However, it is by no means clear that} \he{.hEsEd} \hl{necessarily involves a covenant or means fidelity to a covenant. Stoebe argues that it refers to an attitude as well as to actions. This attitude is parallel to love, }\he{rA.hUM} \hl{goodness,} \he{.tOb} \hl{etc. lt is a kind of love, including mercy, }\he{.hAn*UN}\hl{, when the object is in a pitiful state. It often takes verbs of action, ``do,'' ``keep,'' and so refers to acts of love as well as to the attribute. The word ``lovingkindness'' of the KJV is archaic, but not far from the fulness of meaning of the word.}
						\end{compactdesc}
				\end{compactdesc}
				
				% \begin{compactdesc}
					% \item[Some OT uses] \textbf{}
						% \begin{compactdesc}
							% \item[]
						% \end{compactdesc}
				% \end{compactdesc}
			
		% % --
		% \paragraph{Septuagint}
		% % --
			% \label{MERCY-LXX}
		
			% %
			% \subparagraph{\texorpdfstring{\gr{sample word}}{}}
			% %
		
	% % -----
	% \subsubsection{New Testament} % *NT
	% % -----
		% \label{MERCY-NT}
	
		% % --
		% \paragraph{KJV}
		% % --
			% \label{MERCY-NT-KJV}
			
	
			% \begin{tabular}{ll}
				% Total occurrences in the NT & 
			% \end{tabular}
			
			% \begin{compactdesc}
				% \item[]
			% \end{compactdesc}
			
		% % --
		% \paragraph{Greek New Testament} % *GREEK
		% % --
			% \label{MERCY-GREEK}
		
			% %
			% \subparagraph{\texorpdfstring{\gr{sample word}}{}}
			% %
				% \label{KATALLAGE}
			
				% \begin{compactdesc}
					% \item[BDAG , PDF ] \textbf{}
						% \begin{compactitem}
							% \item 
						% \end{compactitem}
					% \item[LSJ , PDF ] \textbf{}
						% \begin{compactitem}
							% \item 
						% \end{compactitem}
				% \end{compactdesc}
				
				% \begin{compactdesc}
					% \item[Some NT uses] \textbf{}
						% \begin{compactdesc}
							% \item[]
						% \end{compactdesc}
				% \end{compactdesc}
		
	% % -----
	% \subsubsection{Book of Mormon} % *BOM
	% % -----
		% \label{MERCY-BOM}
	
		% \begin{tabular}{ll}
			% Total occurrences in the BOM & 
		% \end{tabular}
		
		% \begin{compactdesc}
			% \item[]
		% \end{compactdesc}
		
	% -----
	\subsubsection{Doctrine and Covenants} % *DC
	% -----
		\label{MERCY-DC}
	
		% \begin{tabular}{ll}
			% Total occurrences in the DC & 
		% \end{tabular}
		
		\begin{compactdesc}
			\item[{\hyperref[DC97-2]{DC 97:2}}]
				\begin{compactitem}
					\item The whole first half of this section is a good place to look at how the Lord extends mercy.
				\end{compactitem}
		\end{compactdesc}
		
	% % -----
	% \subsubsection{Pearl of Great Price} % *PGP
	% % -----
		% \label{MERCY-PGP}
	
		% \begin{tabular}{ll}
			% Total occurrences in the PGP & 
		% \end{tabular}
		
		% \begin{compactdesc}
			% \item[]
		% \end{compactdesc}
		
% ----------------------
\subsection{Key Phrases} % *KEYPHRASES *PHRASES
% ----------------------
	\label{MERCY-PHRASES}

	\begin{compactdesc}
	
		\item[arm of mercy] \textbf{7} | \hyperref[JACOB-6-5]{Jacob 6:5}; \hyperref[MOSIAH-16-12]{Mosiah 16:12} (2); \hyperref[MOSIAH-29-20]{Mosiah 29:20}; \hyperref[ALMA-5-33]{Alma 5:33}; \hyperref[3NEPHI-9-14]{3 Nephi 9:14}; \hyperref[DC29-1]{DC 29:1}
		
		\item[right(s) of mercy] \textbf{1} | Moroni 7:27
			\begin{compactdesc}
				\item[Commentary] See the definitions of \hyperref[RIGHT-NOUN-1828]{right} in the 1828.
			\end{compactdesc}
			
		\item[tender mercy] \textbf{16} | \hyperref[PSALMS-25-6]{Psalms 25:6}; \hyperref[PSALMS-40-11]{Psalms 40:11}; \hyperref[PSALMS-51-1]{Psalms 51:1}; \hyperref[PSALMS-69-16]{Psalms 69:16}; \hyperref[PSALMS-77-9]{Psalms 77:9}; \hyperref[PSALMS-79-8]{Psalms 79:8}; \hyperref[PSALMS-103-4]{Psalms 103:4}; \hyperref[PSALMS-119-77]{Psalms 119:77}; \hyperref[PSALMS-119-156]{Psalms 119:156}; \hyperref[PSALMS-145-9]{Psalms 145:9}; \hyperref[PROVERBS-12-10]{Proverbs 12:10}; \hyperref[LUKE-1-78]{Luke 1:78}; \hyperref[JAMES-5-11]{James 5:11}; \hyperref[1NEPHI-1-20]{1 Nephi 1:20}; \hyperref[1NEPHI-8-8]{1 Nephi 8:8}; \hyperref[ETHER-6-12]{Ether 6:12}
			\begin{compactdesc}
				\item[Bible anchor verses] \phantom{.}
					\begin{compactdesc}
						\item[Psalms 25:6] ...\he{ra.ha:mEyKA}
						\item[Psalms 103:4] ...\he{w:ra.ha:miyM}
						\item[Proverbs 12:10] ...\he{w:ra.ha:mey}
						\item[Luke 1:78] ...\gr{σπλάγχνα ἐλέους}
						\item[James 5:11] ...\gr{οἰκτίρμων}
					\end{compactdesc}
				\item[Commentary] ---
			\end{compactdesc}
		
	\end{compactdesc}
	
% % ----------------------
% \subsection{Bible Dictionaries} % *DICTIONARIES
% % ----------------------
	% \label{MERCY-BD}

% % ----------------------
% \subsection{Restoration Usage} % *RESTORATION
% % ----------------------
	% \label{MERCY-RESTORATION}

	% % -----
	% \subsubsection{Joseph Smith} % *JS
	% % -----
		% \label{MERCY-JS}
	
		% \begin{compactdesc}
			% \item[DATE/SOURCE]
		% \end{compactdesc}
	
	% % -----
	% \subsubsection{Brigham Young} % *BY
	% % -----
		% \label{MERCY-BY}
	
		% \begin{compactdesc}
			% \item[DATE/SOURCE]
		% \end{compactdesc}
	
% % ----------------------
% \subsection{Commentary and Studies} % *COMMENTARY *STUDIES
% % ----------------------
	% \label{MERCY-COMMENTARY}

	% % -----
	% \subsubsection{Key Points}
	% % -----
		% \label{MERCY-KEYS}
	
		% \begin{compactitem}
			% \item
		% \end{compactitem}
		
	% % -----
	% \subsubsection{Sample Study}
	% % -----
		% \label{MERCY-study}